% chapters/13/latihan.tex
\bagian{Latihan dan Refleksi}

\begin{enumerate}
  \item Jelaskan apa yang dimaksud dengan \istilah{greyboxing}, dan sebutkan
        dua alasan mengapa level pertama sebaiknya disusun dengan cara itu.

  \item Sebutkan tiga ukuran level yang perlu ditetapkan lebih dahulu, lalu
        jelaskan mengapa ketiganya saling bergantung.

  \item Jelaskan perbedaan \istilah{Collider} biasa dan \istilah{trigger},
        beserta satu contoh pemakaian masing-masing.

  \item Sebuah benda hendak dibuat menjadi pintu yang membuka ketika pemain
        mendekat. Tipe \istilah{Rigidbody} mana yang tepat, dan mengapa?

  \item Jelaskan mengapa benda diam tidak perlu diberi \istilah{Rigidbody}.
        Apa yang terjadi apabila tetap diberikan?

  \item Jelaskan mengapa masukan pemain dibaca pada \code{Update} sedangkan
        gerak \istilah{Rigidbody} dikerjakan pada \code{FixedUpdate}.

  \item Pemain bergetar ketika berjalan. Sebutkan dua penyebab yang mungkin,
        lalu jelaskan cara memeriksa masing-masing.

  \item Sebuah benda bergerak 30 satuan per detik dan lantainya tebal 0,1
        satuan. Tentukan apakah benda itu dapat menembus lantai pada langkah
        fisika 0,02 detik, beserta perhitungannya, lalu sebutkan dua cara
        perbaikannya.
\end{enumerate}

\begin{refleksibox}
\begin{enumerate}
  \item Pada langkah kelima praktikum, masalah fisika apa yang Anda temui?
        Berapa lama Anda mencarinya, dan pada akhirnya apakah penyebabnya
        berada pada kode atau pada pengaturan komponen?

  \item Ketika Anda menambahkan dengan sengaja salah satu masalah dari
        Tabel~\ref{tab:masalah-fisika}, apakah gejalanya sesuai dengan yang
        tertulis pada tabel? Apabila tidak, apa yang berbeda?

  \item Bandingkan pengalaman Anda menyusun level pada Unity dengan pengalaman
        Anda menyusun peta pada Bab~X. Bagian mana yang lebih mudah, dan
        mengapa?
\end{enumerate}
\end{refleksibox}
